\section{Requerimientos Previos y Entorno de Trabajo}

La correcta implementación de un clúster Hadoop mediante contenedores Docker requiere disponer de un entorno técnico adecuado y una estructura de proyecto organizada que facilite la construcción, despliegue y administración de los servicios distribuidos. En esta sección se presentan los requerimientos previos, las características del equipo utilizado, el software necesario para la ejecución del laboratorio, así como la descripción detallada de los archivos que componen la solución.

\subsection{Especificaciones del Computador}

Para garantizar un desempeño adecuado durante la construcción de imágenes, la inicialización de contenedores y la ejecución de procesos distribuidos, se recomienda utilizar un equipo con características mínimas que permitan la asignación suficiente de recursos a los servicios del clúster Hadoop. La Tabla~\ref{tab:especificaciones} resume las especificaciones recomendadas y las utilizadas durante el desarrollo del laboratorio.

\begin{table}[H]
\centering
\begin{tabular}{|l|p{9cm}|}
\hline
\textbf{Componente} & \textbf{Especificación Recomendada / Utilizada} \\ \hline
Procesador & CPU multinúcleo (mínimo 4 núcleos físicos). \\ \hline
Memoria RAM & 8 GB mínimo (16 GB recomendado para mejor desempeño). \\ \hline
Almacenamiento & 20 GB libres como mínimo para imágenes, contenedores y archivos HDFS. \\ \hline
Sistema Operativo & Windows 10/11, Linux o macOS con soporte para Docker. \\ \hline
Virtualización & Virtualización por hardware habilitada (Intel VT-x / AMD-V). \\ \hline
\end{tabular}
\caption{Especificaciones del equipo para el laboratorio.}
\label{tab:especificaciones}
\end{table}

\subsection{Software Requerido}

Para la ejecución del clúster Hadoop basado en Docker, se requiere el siguiente software:

\begin{itemize}
    \item \textbf{Docker Engine}: Motor de contenedores encargado de construir imágenes, crear contenedores y gestionar redes.
    \item \textbf{Docker Compose}: Herramienta para definir e implementar aplicaciones multicontenedor mediante archivos YAML.
    \item \textbf{OpenJDK (incluido en la imagen personalizada)}: Dependencia esencial para ejecutar los servicios de Hadoop.
    \item \textbf{Sistema de archivos Unix-like (dentro del contenedor)}: Necesario para la ejecución de scripts \texttt{.sh} y gestión de demonios.
\end{itemize}

La combinación de estas herramientas permite construir un entorno aislado, reproducible y portable para desplegar el ecosistema Hadoop.

\subsection{Docker como Plataforma de Ejecución}

Docker actúa como una plataforma ligera que abstrae las dependencias del sistema anfitrión y permite que Hadoop se ejecute dentro de contenedores independientes. Su utilización en este laboratorio facilita:

\begin{itemize}
    \item La rápida creación y destrucción de nodos.
    \item La estandarización de configuraciones.
    \item La replicabilidad del entorno sin necesidad de máquinas virtuales pesadas.
    \item El aislamiento entre servicios críticos como NameNode, DataNode, ResourceManager y NodeManager.
\end{itemize}

\subsection{Imágenes Utilizadas}

El laboratorio emplea una imagen de Hadoop personalizada construida mediante un archivo \texttt{Dockerfile}. Esta imagen contiene:

\begin{itemize}
    \item OpenJDK 8 como entorno de ejecución.
    \item Variables de entorno para Hadoop.
    \item Copia de archivos de configuración desde el host hacia el contenedor.
    \item Script de inicialización del clúster.
\end{itemize}

Adicionalmente, Docker Compose se encarga de generar contenedores independientes para:

\begin{itemize}
    \item \textbf{NameNode}: Nodo maestro encargado de la gestión de metadatos del HDFS.
    \item \textbf{DataNode(s)}: Nodos de almacenamiento.
    \item \textbf{ResourceManager}: Administrador general de recursos YARN.
    \item \textbf{NodeManager(s)}: Ejecutores de tareas en los nodos.
\end{itemize}

\subsection{Estructura del Proyecto}

El proyecto implementado posee la siguiente estructura de carpetas, organizada con el objetivo de separar claramente los archivos de configuración, los scripts de inicialización y los recursos necesarios para el despliegue del clúster:

    \begin{verbatim}
    Hadoop/
    |
    +-- hadoop-conf/
    |   +-- core-site.xml
    |   +-- hadoop-env.sh
    |   +-- hdfs-site.xml
    |   +-- log4j.properties
    |   +-- mapred-site.xml
    |   +-- capacity-scheduler.xml
    |   +-- workers
    |   +-- yarn-site.xml
    |
    +-- docker-compose.yml
    +-- Dockerfile
    +-- start-hadoop.sh
    \end{verbatim}

A continuación se detalla la función de cada uno de estos elementos.

\subsection{Descripción de Archivos y Directorios}

\subsubsection{Directorio \texttt{hadoop-conf/}}

Este directorio contiene todos los archivos de configuración esenciales para la correcta ejecución de HDFS, YARN y MapReduce. Entre ellos se encuentran:

\begin{itemize}
    \item \textbf{core-site.xml}: Define parámetros globales del sistema, incluyendo la URI del NameNode.
    \item \textbf{hadoop-env.sh}: Contiene variables de entorno internas de Hadoop, como rutas del \texttt{JAVA\_HOME}.
    \item \textbf{hdfs-site.xml}: Configura la replicación, directorios de almacenamiento y comportamiento del namenode/datanode.
    \item \textbf{log4j.properties}: Define niveles de registro y salida de logs.
    \item \textbf{mapred-site.xml}: Parámetros específicos del framework MapReduce.
    \item \textbf{capacity-scheduler.xml}: Configuración avanzada para asignación de recursos YARN (opcional en este laboratorio).
    \item \textbf{workers}: Lista de nodos trabajadores (DataNodes y NodeManagers).
    \item \textbf{yarn-site.xml}: Configuración del sistema de gestión de recursos YARN.
\end{itemize}

\subsubsection{\texttt{Dockerfile}}

Archivo utilizado para construir la imagen base de Hadoop. Incluye:

\begin{itemize}
    \item Instalación de dependencias.
    \item Configuración de variables de entorno.
    \item Copia de archivos de configuración hacia el contenedor.
    \item Ejecución del script de inicio del clúster.
\end{itemize}

\subsubsection{\texttt{docker-compose.yml}}

Archivo central para el despliegue de la arquitectura distribuida. Define:

\begin{itemize}
    \item Servicios del NameNode, DataNodes, ResourceManager y NodeManagers.
    \item Redes virtuales internas.
    \item Volúmenes persistentes para el almacenamiento del HDFS.
    \item Dependencias y orden de arranque.
\end{itemize}

\subsubsection{\texttt{start-hadoop.sh}}

Script encargado de inicializar todos los servicios de Hadoop, incluyendo:

\begin{itemize}
    \item Formateo del NameNode (solo si es necesario).
    \item Inicio de los servicios HDFS.
    \item Inicio de los servicios YARN.
    \item Verificación del estado general del clúster.
\end{itemize}

Este archivo automatiza el proceso completo de puesta en marcha del sistema distribuido.
